
\subsection{Comparison with State-of-the-art}
\label{Compare_SOTA}
\begin{table}[h]
\centering
\caption{Results of accuracy and fairness of different methods on Fitzpatrick-17k dataset, using skin tone as the sensitive attribute. The dark skin is the privileged group with higher accuracy by vanilla training. 
The results of \emph{VGG-11}, \emph{AdvConf}, \emph{AdvRev}, \emph{DomainIndep} and \emph{FairPrune} are the experimental results reported in \cite{wu2022fairprune}.
}
\setlength{\tabcolsep}{6pt}
\label{Table:compare_fitz}
\scriptsize
\begin{tabular}{c c c c c c c c}\toprule
% \hline
\multicolumn{1}{c}{\text{}} & \multicolumn{1}{c }{\text{}} &\multicolumn{3}{c }{\text{Accuracy}} & \multicolumn{3}{c }{\text{Fairness}}  \\ 
\cmidrule(lr){3-5}
\cmidrule(ll){6-8}
\multicolumn{1}{ c }{\text{Method}} & \text{Skin Tone} &\text{Precision} & \text{Recall} & \text{F1-score} & \text{Eopp0 $\downarrow$} & \text{Eopp1 $\downarrow$} & \text{Eodd $\downarrow$} \\ 
% \cmidrule(lr){1-1}
% \cmidrule(lr){2-2}
% \cmidrule(lr){3-5}
% \cmidrule(ll){6-8}
\hline
\hline
\multirow{4}{*}{VGG-11} & \text{Dark} & 0.563 & 0.581 & 0.546 & \multirow{4}{*}{0.0013} & \multirow{4}{*}{0.361} & \multirow{4}{*}{0.182} \\ 
 & \text{Light} & 0.482 & 0.495 & 0.473 &  &  & \\ 
 & \text{Avg. $\uparrow$} & 0.523 & 0.538 & 0.510 &  &  & \\ 
 & \text{Diff. $\downarrow$} & 0.081 & 0.086 & 0.073 &  &  & \\ 
\hline
\multicolumn{1}{ c }{\multirow{4}{*}{AdvConf}} & \text{Dark} & 0.506 & 0.562 & 0.506 & \multirow{4}{*}{0.0011} & \multirow{4}{*}{0.339} & \multirow{4}{*}{0.169} \\ 
 & \text{Light} & 0.427 & 0.464 & 0.426 &  &  & \\ 
 & \text{Avg. $\uparrow$} & 0.467 & 0.513 & 0.466 &  &  & \\ 
 & \text{Diff. $\downarrow$} & 0.079 & 0.098 & 0.080 &  &  & \\ 
\hline
\multicolumn{1}{ c }{\multirow{4}{*}{AdvRev}} & \text{Dark} & 0.514 & 0.545 & 0.503 & \multirow{4}{*}{0.0011} & \multirow{4}{*}{0.334} & \multirow{4}{*}{0.166} \\ 
 & \text{Light} & 0.489 & 0.469 & 0.457 &  &  & \\ 
 & \text{Avg. $\uparrow$} & 0.502 & 0.507 & 0.480 &  &  & \\ 
 & \text{Diff. $\downarrow$} & 0.025 & 0.076 & 0.046 &  &  & \\ 
\hline
\multicolumn{1}{ c }{\multirow{4}{*}{DomainIndep}} & \text{Dark} & 0.559 & 0.540 & 0.530 & \multirow{4}{*}{0.0012} & \multirow{4}{*}{0.323} & \multirow{4}{*}{0.161} \\ 
 & \text{Light} & 0.541 & 0.529 & 0.512 &  &  & \\ 
 & \text{Avg. $\uparrow$} & 0.550 & 0.534 & 0.521 &  &  & \\ 
 & \text{Diff. $\downarrow$} & 0.018 & 0.010 & 0.018 &  &  & \\ 
\hline
\multicolumn{1}{ c }{\multirow{4}{*}{HSIC}} & \text{Dark} & 0.548 & 0.522 & 0.513 & \multirow{4}{*}{0.0013} & \multirow{4}{*}{0.331} & \multirow{4}{*}{0.166} \\ 
 & \text{Light} & 0.513 & 0.506 & 0.486 &  &  & \\ 
 & \text{Avg.$\uparrow$} & 0.530 & 0.515 & 0.500 &  &  & \\ 
 & \text{Diff. $\downarrow$} & 0.040 & 0.018 & 0.029 &  &  & \\ 
\hline

% \multicolumn{1}{ c }{HSIC} & 0.74 & 15 & 630 & 0.70 & 15 & 618 & -\\ 
\multicolumn{1}{ c }{\multirow{4}{*}{MFD}} & \text{Dark} & 0.514 & 0.545 & 0.503 & \multirow{4}{*}{0.0011} & \multirow{4}{*}{0.334} & \multirow{4}{*}{0.166} \\ 
 & \text{Light} & 0.489 & 0.469 & 0.457 &  &  & \\ 
 & \text{Avg. $\uparrow$} & 0.502 & 0.507 & 0.480 &  &  & \\ 
 & \text{Diff. $\downarrow$} & 0.025 & 0.076 & 0.046 &  &  & \\ 
\hline
\hline
% \multicolumn{1}{ c }{MFD} & 0.06 & 7 & 21 & 0.06 & 5 & 27 & -\\ 
% \multirow{4}{*}{\multicolumn{1}{ c }{DomainIndep}} & \text{Dark} & 0.559 & 0.540 & 0.530 & \multirow{4}{*}{0.0012} & \multirow{4}{*}{0.323} & \multirow{4}{*}{0.161} \\ 
%  & \text{Light} & 0.541 & 0.529 & 0.512 &  &  & \\ 
%  & \text{Avg.} & 0.550 & 0.534 & 0.521 &  &  & \\ 
%  & \text{Diff.} & 0.018 & 0.010 & 0.018 &  &  & \\ 
% \hline
% \multicolumn{1}{ c }{DomainIndep} & 0.06 & 7 & 21 & 0.06 & 5 & 27 & -\\ 
{\multirow{4}{*}{FairPrune}} & \text{Dark} & 0.567 & 0.519 & 0.507 & \multirow{4}{*}{\textbf{0.0008}} & \multirow{4}{*}{0.330} & \multirow{4}{*}{0.165} \\ 
 & \text{Light} & 0.496 & 0.477 & 0.459 &  &  & \\ 
 & \text{Avg. $\uparrow$} & 0.531 & 0.498 & 0.483 &  &  & \\ 
 & \text{Diff. $\downarrow$} & 0.071 & 0.042 & 0.048 &  &  & \\ 
\hline
% \multicolumn{1}{ c }{FairPrune} & 0.06 & 7 & 21 & 0.06 & 5 & 27 & -\\ 
\multicolumn{1}{ c }{\multirow{4}{*}{ME-FairPrune}} & \text{Dark} & 0.564 & 0.529 & 0.523 & \multirow{4}{*}{0.0012} & \multirow{4}{*}{\textbf{0.305}} & \multirow{4}{*}{\textbf{0.152}} \\ 
 & \text{Light} & 0.542 & 0.535 & 0.522 &  &  & \\ 
 & \text{Avg. $\uparrow$} & 0.553 & 0.532 & 0.522 &  &  & \\ 
 & \text{Diff. $\downarrow$} & 0.022 & 0.006 & 0.001 &  &  & \\
 % \hline
\bottomrule
\end{tabular}
\end{table}


% \begin{tabular}{SSSSSSSS} 
% \toprule
% \multicolumn{7}{c}{Accuracy} {Fairness}\\
% \cmidrule(r){3-5} \cmidrule(lr){6-8}
%     {Method} & {Skin Tone} & {Precision} & {Recall} & {F1-score} & {Eopp0} & {Eopp1} & {Eodd} \\ \midrule
%     1  & 16.128 & +8.872 & 16.128 & 1.402 & 1.373 & -146.6 & -137.6 \\
%     2  & 3.442  & -2.509 & 3.442  & 0.299 & 0.343 & 133.2  & 152.4  \\
%     3  & 1.826  & -0.363 & 1.826  & 0.159 & 0.119 & 168.5  & -161.1 \\
%     4  & 0.993  & -0.429 & 0.993  & 0.086 & 0.08  & 25.6   & 90     \\ \midrule
%     5  & 1.29   & +0.099 & 1.29   & 0.112 & 0.097 & -175.6 & -114.7 \\
%     6  & 0.483  & -0.183 & 0.483  & 0.042 & 0.063 & 22.3   & 122.5  \\
%     7  & 0.766  & -0.475 & 0.766  & 0.067 & 0.039 & 141.6  & -122   \\
%     8  & 0.624  & +0.365 & 0.624  & 0.054 & 0.04  & -35.7  & 90     \\ \midrule
%     9  & 0.641  & -0.466 & 0.641  & 0.056 & 0.045 & 133.3  & -106.3 \\
%     10 & 0.45   & +0.421 & 0.45   & 0.039 & 0.034 & -69.4  & 110.9  \\
%     11 & 0.598  & -0.597 & 0.598  & 0.052 & 0.025 & 92.3   & -109.3 \\ \bottomrule
% \end{tabular}

% In this section, we applied our framework to the current state-of-the-art method FairPrune \cite{wu2022fairprune}, and compare the results with several baselines, including CNN (ResNet18 and VGG-11), AdvConf \cite{alvi2018turning}, AdvRev \cite{zhang2018mitigating}, DomainIndep \cite{wang2020towards}, HSIC \cite{quadrianto2019discovering}, and MFD \cite{jung2021fair}. For each dataset, we report accuracy and fairness results, including precision, recall, and F1-score metrics.

In this section, we compare our framework with several baselines, including CNN (ResNet18 and VGG-11), AdvConf \cite{alvi2018turning}, AdvRev \cite{zhang2018mitigating}, DomainIndep \cite{wang2020towards}, HSIC \cite{quadrianto2019discovering}, and MFD \cite{jung2021fair}. We also compare our framework to the current state-of-the-art method FairPrune \cite{wu2022fairprune}. For each dataset, we report accuracy and fairness results, including precision, recall, and F1-score metrics.

\textbf{ISIC 2019 Dataset.}
In Table \ref{Table:compare_isic}, ME-FairPrune refers to the FairPrune applied in our framework. It shows that our ME framework has improved all fairness scores and accuracy in all average scores when applied to FairPrune. Additionally, the difference in each accuracy metric is smaller than that of the original FairPrune. This is because the ME framework ensures that the early exited instances have a high level of confidence in their correctness, and the classification through shallower, fairer features further improves fairness. 

\textbf{Fitzpatrick-17k Dataset.}
To evaluate the extensibility of our framework on different model structures, we use VGG-11 as the backbone of the Fitzpatrick-17k dataset. Table \ref{Table:compare_fitz} demonstrates that the ME framework outperforms all other methods with the best Eopp1 and Eodd scores, showing a 7.5\% and 7.9\% improvement over FairPrune, respectively. Similar to the ISIC 2019 dataset, our results show better mean accuracy and more minor accuracy differences in all criteria than FairPrune. Furthermore, the F1-score and Precision average value are superior to other methods, which show 8.1\% and 4.1\% improvement over FairPrune, respectively.

\begin{table}[h]
\centering
\caption{Accuracy and fairness of classification results across different baselines with and without the ME training framework on the ISIC 2019 dataset.}
\setlength{\tabcolsep}{6pt}
\label{Table:compare_isic_me}
\scriptsize
\begin{tabular}{c c c c c c c c}\toprule
% \hline
\multicolumn{1}{c }{\text{}} & \multicolumn{1}{c }{\text{}} &\multicolumn{3}{c }{\text{Accuracy}} & \multicolumn{3}{c }{\text{Fairness}}  \\ 
\cmidrule(lr){3-5}
\cmidrule(ll){6-8}
\multicolumn{1}{ c }{\text{Method}} & \text{Gender} &\text{Precision} & \text{Recall} & \text{F1-score} & \text{Eopp0 $\downarrow$} & \text{Eopp1 $\downarrow$} & \text{Eodd $\downarrow$} \\ 
\hline
\hline

\multicolumn{1}{ c }{\multirow{4}{*}{ResNet18}} & \text{Female} & 0.793 & 0.721 & 0.746 & \multirow{4}{*}{\textbf{0.006}} & \multirow{4}{*}{0.044} & \multirow{4}{*}{0.022} \\ 
 & \text{Male} & 0.731 & 0.725 & 0.723 &  &  & \\ 
 & \text{Avg. $\uparrow$} & 0.762 & 0.723 & 0.735 &  &  & \\ 
 & \text{Diff. $\downarrow$} & 0.063 & 0.004 & 0.023 &  &  & \\ 
\hline
% \cmidrule(lr){1-1}
% \cmidrule(lr){2-2}
% \cmidrule(lr){3-5}
% \cmidrule(ll){6-8}

\multicolumn{1}{ c }{\multirow{4}{*}{ME-ResNet18}} & \text{Female} & 0.748 & 0.724 & 0.733 & \multirow{4}{*}{\textbf{0.006}} & \multirow{4}{*}{\textbf{0.031}} & \multirow{4}{*}{\textbf{0.016}} \\ 
 & \text{Male} & 0.723 & 0.736 & 0.726 &  &  & \\ 
 & \text{Avg. $\uparrow$} & 0.735 & 0.730 & 0.730 &  &  & \\ 
 & \text{Diff. $\downarrow$} & 0.025 & 0.012 & 0.007 &  &  & \\ 
\hline
\hline

\multicolumn{1}{ c }{\multirow{4}{*}{HSIC}} & \text{Female} & 0.744 & 0.660 & 0.696 & \multirow{4}{*}{0.008} & \multirow{4}{*}{0.042} & \multirow{4}{*}{0.020} \\ 
 & \text{Male} & 0.718 & 0.697 & 0.705 &  &  & \\ 
 & \text{Avg. $\uparrow$} & 0.731 & 0.679 & 0.700 &  &  & \\ 
 & \text{Diff. $\downarrow$} & 0.026 & 0.037 & 0.009 &  &  & \\ 
\hline

\multicolumn{1}{ c }{\multirow{4}{*}{ME-HSIC}} & \text{Female} & 0.733 & 0.707 & 0.716 & \multirow{4}{*}{\textbf{0.007}} & \multirow{4}{*}{{\textbf{0.034}}} & \multirow{4}{*}{{\textbf{0.018}}} \\ 
 & \text{Male} & 0.713 & 0.707 & 0.707 &  &  & \\ 
 & \text{Avg. $\uparrow$} & 0.723 & 0.707 & 0.712 &  &  & \\ 
 & \text{Diff. $\downarrow$} & 0.020 & 0.000 & 0.009 &  &  & \\ 
\hline
\hline

% \multicolumn{1}{ c }{HSIC} & 0.74 & 15 & 630 & 0.70 & 15 & 618 & -\\ 
\multicolumn{1}{ c }{\multirow{4}{*}{MFD}} & \text{Female} & 0.770 & 0.697 & 0.726 & \multirow{4}{*}{\textbf{0.005}} & \multirow{4}{*}{0.051} & \multirow{4}{*}{0.024} \\ 
 & \text{Male} & 0.772 & 0.726 & 0.744 &  &  & \\ 
 & \text{Avg. $\uparrow$} & 0.771 & 0.712 & 0.735 &  &  & \\ 
 & \text{Diff. $\downarrow$} & 0.002 & 0.029 & 0.018 &  &  & \\ 
\hline

\multicolumn{1}{ c }{\multirow{4}{*}{ME-MFD}} & \text{Female} & 0.733 & 0.698 & 0.711 & \multirow{4}{*}{\textbf{0.005}} & \multirow{4}{*}{\textbf{0.024}} & \multirow{4}{*}{\textbf{0.012}} \\ 
 & \text{Male} & 0.772 & 0.713 & 0.739 &  &  & \\ 
 & \text{Avg. $\uparrow$} & 0.752 & 0.706 & 0.725 &  &  & \\ 
 & \text{Diff. $\downarrow$} & 0.039 & 0.015 & 0.028 &  &  & \\ 

% \multicolumn{1}{ c }{MFD} & 0.06 & 7 & 21 & 0.06 & 5 & 27 & -\\ 
% \multirow{4}{*}{\multicolumn{1}{ c }{DomainIndep}} & \text{Female} & 0.729 & 0.747 & 0.734 & \multirow{4}{*}{0.010} & \multirow{4}{*}{0.086} & \multirow{4}{*}{0.042} \\ 
%  & \text{Male} & 0.725 & 0.694 & 0.702 &  &  & \\ 
%  & \text{Avg.} & 0.727 & 0.721 & 0.718 &  &  & \\ 
%  & \text{Diff.} & 0.004 & 0.053 & 0.031 &  &  & \\ 
% \hline
% \multicolumn{1}{ c }{DomainIndep} & 0.06 & 7 & 21 & 0.06 & 5 & 27 & -\\ 
% \hline
\bottomrule
\end{tabular}
\end{table}


% \begin{tabular}{SSSSSSSS} 
% \toprule
% \multicolumn{7}{c}{Accuracy} {Fairness}\\
% \cmidrule(r){3-5} \cmidrule(lr){6-8}
%     {Method} & {Skin Tone} & {Precision} & {Recall} & {F1-score} & {Eopp0} & {Eopp1} & {Eodd} \\ \midrule
%     1  & 16.128 & +8.872 & 16.128 & 1.402 & 1.373 & -146.6 & -137.6 \\
%     2  & 3.442  & -2.509 & 3.442  & 0.299 & 0.343 & 133.2  & 152.4  \\
%     3  & 1.826  & -0.363 & 1.826  & 0.159 & 0.119 & 168.5  & -161.1 \\
%     4  & 0.993  & -0.429 & 0.993  & 0.086 & 0.08  & 25.6   & 90     \\ \midrule
%     5  & 1.29   & +0.099 & 1.29   & 0.112 & 0.097 & -175.6 & -114.7 \\
%     6  & 0.483  & -0.183 & 0.483  & 0.042 & 0.063 & 22.3   & 122.5  \\
%     7  & 0.766  & -0.475 & 0.766  & 0.067 & 0.039 & 141.6  & -122   \\
%     8  & 0.624  & +0.365 & 0.624  & 0.054 & 0.04  & -35.7  & 90     \\ \midrule
%     9  & 0.641  & -0.466 & 0.641  & 0.056 & 0.045 & 133.3  & -106.3 \\
%     10 & 0.45   & +0.421 & 0.45   & 0.039 & 0.034 & -69.4  & 110.9  \\
%     11 & 0.598  & -0.597 & 0.598  & 0.052 & 0.025 & 92.3   & -109.3 \\ \bottomrule
% \end{tabular}

\subsection{Multi-Exit Training on Different Method}
\label{ME_different_baseline}
In this section, we evaluate the performance of our ME framework when applied to different methods. Table \ref{Table:compare_isic_me} presents the results of our experiments on ResNet18, MFD, and HSIC. Our ME framework improved the Eopp1 and Eodd scores of the original ResNet18 model, which did not apply fairness regularization loss, $l_s$, in total loss. Furthermore, our framework achieved comparable performance in terms of fairness to FairPrune, as shown in Table \ref{Table:compare_isic}. This demonstrates its potential to achieve fairness without using sensitive attributes during training.

We also applied our ME framework to MFD and HSIC, which initially exhibited better fairness performance than other baselines. With our framework, these models showed better fairness while maintaining similar levels of accuracy. These findings suggest that our ME framework can improve the fairness of existing models, making them more equitable without compromising accuracy.
% of accuracy
% \begin{figure}
% % \setlength{\belowcaptionskip}{-20pt} 
% \begin{center}

% \includegraphics[width=0.95\linewidth]{Fig/Ablation.png}
% % \includegraphics[width=0.9\columnwidth]{Fig/Method.png}
% \caption{Ablation study on (a) the early exit confidence threshold $\theta$ and (b) the early exit strategy for gender fairness on ISIC 2019 dataset. The $CLF_{1}$ through $CLF_{4}$ refer to internal classifiers, while $CLF_{f}$ refers to the final classifier.}
% \label{Fig.Ablation}
% \end{center}
% \end{figure}

\subsection{Ablation Study}

% \begin{figure}[h]
% % \setlength{\belowcaptionskip}{-20pt} 
% \begin{center}

% \includegraphics[width=0.95\linewidth]{Fig/Ablation.png}
% % \includegraphics[width=0.9\columnwidth]{Fig/Method.png}
% \caption{Ablation study on early exit threshold $\theta$ and the inference strategy for gender fairness on ISIC 2019 dataset.}
% \label{Fig.Ablation}
% \end{center}
% \end{figure}

% In this section, we conduct an ablation study to investigate the effectiveness of different inference thresholds and the impact of early exits on the ME-FairPrune method. 

\begin{figure}
% \setlength{\belowcaptionskip}{-20pt} 
\begin{center}

\includegraphics[width=0.93\linewidth]{Fig/Ablation.png}
% \includegraphics[width=0.9\columnwidth]{Fig/Method.png}
\caption{Ablation study on (a) the early exit confidence threshold $\theta$ and (b) the early exit strategy for gender fairness on ISIC 2019 dataset. The $CLF_{1}$ through $CLF_{4}$ refer to internal classifiers, while $CLF_{f}$ refers to the final classifier.}
\label{Fig.Ablation}
\vspace{-20pt}
\end{center}
\end{figure}

\textbf{\indent Effect of Different Confidence Thresholds.} To investigate the impact of varying confidence thresholds $\theta$ on accuracy and fairness, we apply the ME-FairPrune method to a pre-trained model from the ISIC 2019 dataset and test different thresholds. Our results, shown in Fig. \ref{Fig.Ablation} (a), indicate that increasing the threshold improves accuracy and fairness. Thus, we recommend setting the threshold to 0.999 for optimal performance.

\textbf{Effect of Early Exits.} In Fig. \ref{Fig.Ablation} (b), we compare ME-FairPrune using an early exit policy with exiting from each specific exit. The results show that across all criteria, no specific $CLF$ outperforms the early exit policy in terms of Eodd, while our early exit policy achieves an accuracy level comparable to the original classifier output $CLF_{f}$. These findings underscore the importance of the proposed early exit strategy for achieving optimal performance.






% We compared our models with several baselines. \emph{Vanilla} is a standard CNN without applying a debiased strategy. \emph{AdvConf} \cite{alvi2018turning} and \emph{AdvRev} \cite{zhang2018mitigating} are adversarial training based de-bias method. \emph{HSIC} \cite{quadrianto2019discovering} and \emph{MFD} \cite{jung2021fair} are both regularization-based de-bias method which add the regularization term at loss function to add the fairness constraint in the training process. \emph{DomainIndep} \cite{wang2020towards} uses multiple classifier to explicitly separate information of each group. \emph{FairPrune} \cite{wu2022fairprune} achieves fair dermatology disease classification by pruning the neurons with unfairly favorable to certain group. 

% \subsection{Results on ISIC 2019 Dataset.}

% \begin{table}[ht]
\centering
\caption{Results of accuracy and fairness of different methods on ISIC 2019 dataset, using gender as the sensitive attribute. The female is the privileged group with higher accuracy by vanilla training. All results are our implementation. }
\setlength{\tabcolsep}{6pt}
\label{Table:compare_isic}
\scriptsize
\begin{tabular}{c c c c c c c c}\toprule 
% \hline
% {\multirow{2}{*}{Method}} & \multicolumn{1}{c@{}}{\multirow{2}{*}{Gender}} &\multicolumn{3}{c@{}}{\text{Accuracy}} & \multicolumn{3}{c}{\text{Fairness}}  \\ 
% \cmidrule(lr){3-5}
% \cmidrule(ll){6-8}
% {} & {} &\text{Precision} & \text{Recall} & \text{F1-score} & \text{Eopp0 $\downarrow$} & \text{Eopp1 $\downarrow$} & \text{Eodd $\downarrow$} \\ 
% \cmidrule(lr){1-1}
% \cmidrule(lr){2-2}
% \cmidrule(lr){3-5}
% \cmidrule(ll){6-8}
\multicolumn{1}{c }{\text{}} & \multicolumn{1}{c }{\text{}} &\multicolumn{3}{c }{\text{Accuracy}} & \multicolumn{3}{c }{\text{Fairness}}  \\ 
\cmidrule(lr){3-5}
\cmidrule(ll){6-8}
\multicolumn{1}{ c }{\text{Method}} & \text{Gender} &\text{Precision} & \text{Recall} & \text{F1-score} & \text{Eopp0 $\downarrow$} & \text{Eopp1 $\downarrow$} & \text{Eodd $\downarrow$} \\ 
\hline
\hline
 {\multirow{4}{*}{ResNet18}} & \text{Female} & 0.793 & 0.721 & 0.746 & \multirow{4}{*}{0.006} & \multirow{4}{*}{0.044} & \multirow{4}{*}{0.022} \\ 
 & \text{Male} & 0.731 & 0.725 & 0.723 &  &  & \\ 
 & \text{Avg. $\uparrow$} & 0.762 & 0.723 & 0.735 &  &  & \\ 
 & \text{Diff. $\downarrow$} & 0.063 & 0.004 & 0.023 &  &  & \\ 
% \hline
% \cmidrule(lr){1-1}
% \cmidrule(lr){2-2}
% \cmidrule(lr){3-5}
% \cmidrule(ll){6-8}

% \multirow{4}{*}{\multicolumn{1}{ c }{ME-ResNet18}} & \text{Female} & 0.748 & 0.724 & 0.733 & \multirow{4}{*}{0.006} & \multirow{4}{*}{0.031} & \multirow{4}{*}{0.016} \\ 
%  & \text{Male} & 0.723 & 0.736 & 0.726 &  &  & \\ 
%  & \text{Avg.} & 0.735 & 0.730 & 0.730 &  &  & \\ 
%  & \text{Diff.} & 0.025 & 0.012 & 0.007 &  &  & \\ 
% \hline
\hline


\multirow{4}{*}{AdvConf} & \text{Female} & 0.755 & 0.738 & 0.741 & \multirow{4}{*}{0.008} & \multirow{4}{*}{0.070} & \multirow{4}{*}{0.037} \\ 
 & \text{Male} & 0.710 & 0.757 & 0.731 &  &  & \\ 
 & \text{Avg. $\uparrow$} & 0.733 & 0.747 & 0.736 &  &  & \\ 
 & \text{Diff. $\downarrow$} & 0.045 & 0.020 & 0.010 &  &  & \\ 
\hline
% \cmidrule(lr){1-1}
% \cmidrule(lr){2-2}
% \cmidrule(lr){3-5}
% \cmidrule(ll){6-8}
\multirow{4}{*}{AdvRev} & \text{Female} & 0.778 & 0.683 & 0.716 & \multirow{4}{*}{0.007} & \multirow{4}{*}{0.059} & \multirow{4}{*}{0.033} \\ 
 & \text{Male} & 0.773 & 0.706 & 0.729 &  &  & \\ 
 & \text{Avg. $\uparrow$} & 0.775 & 0.694 & 0.723 &  &  & \\ 
 & \text{Diff. $\downarrow$} & 0.006 & 0.023 & 0.014 &  &  & \\ 
%  \cmidrule(lr){1-1}
% \cmidrule(lr){2-2}
% \cmidrule(lr){3-5}
% \cmidrule(ll){6-8}
\hline
\multirow{4}{*}{DomainIndep} & \text{Female} & 0.729 & 0.747 & 0.734 & \multirow{4}{*}{0.010} & \multirow{4}{*}{0.086} & \multirow{4}{*}{0.042} \\ 
 & \text{Male} & 0.725 & 0.694 & 0.702 &  &  & \\ 
 & \text{Avg. $\uparrow$} & 0.727 & 0.721 & 0.718 &  &  & \\ 
 & \text{Diff. $\downarrow$} & 0.004 & 0.053 & 0.031 &  &  & \\ 
% \hline
\hline
\multirow{4}{*}{HSIC} & \text{Female} & 0.744 & 0.660 & 0.696 & \multirow{4}{*}{0.008} & \multirow{4}{*}{0.042} & \multirow{4}{*}{0.020} \\ 
 & \text{Male} & 0.718 & 0.697 & 0.705 &  &  & \\ 
 & \text{Avg. $\uparrow$} & 0.731 & 0.679 & 0.700 &  &  & \\ 
 & \text{Diff. $\downarrow$} & 0.026 & 0.037 & 0.009 &  &  & \\ 
% \cmidrule(lr){1-1}
% \cmidrule(lr){2-2}
% \cmidrule(lr){3-5}
% \cmidrule(ll){6-8}
% % \hline

% \multirow{4}{*}{\multicolumn{1}{ c }{ME-HSIC}} & \text{Female} & 0.733 & 0.707 & 0.716 & \multirow{4}{*}{0.007} & \multirow{4}{*}{0.034} & \multirow{4}{*}{0.018} \\ 
%  & \text{Male} & 0.713 & 0.707 & 0.707 &  &  & \\ 
%  & \text{Avg.} & 0.723 & 0.707 & 0.712 &  &  & \\ 
%  & \text{Diff.} & 0.020 & 0.000 & 0.009 &  &  & \\ 
% \hline
\hline
% \multicolumn{1}{ c }{HSIC} & 0.74 & 15 & 630 & 0.70 & 15 & 618 & -\\ 
\multicolumn{1}{ c }{\multirow{4}{*}{MFD}} & \text{Female} & 0.770 & 0.697 & 0.726 & \multirow{4}{*}{\textbf{0.005}} & \multirow{4}{*}{0.051} & \multirow{4}{*}{0.024} \\ 
 & \text{Male} & 0.772 & 0.726 & 0.744 &  &  & \\ 
 & \text{Avg. $\uparrow$} & 0.771 & 0.712 & 0.735 &  &  & \\ 
 & \text{Diff. $\downarrow$} & 0.002 & 0.029 & 0.018 &  &  & \\ 
% \hline
% \cmidrule(lr){1-1}
% \cmidrule(lr){2-2}
% \cmidrule(lr){3-5}
% \cmidrule(ll){6-8}

% \multirow{4}{*}{\multicolumn{1}{ c }{ME-MFD}} & \text{Female} & 0.733 & 0.698 & 0.711 & \multirow{4}{*}{0.006} & \multirow{4}{*}{0.024} & \multirow{4}{*}{0.012} \\ 
%  & \text{Male} & 0.772 & 0.713 & 0.739 &  &  & \\ 
%  & \text{Avg.} & 0.752 & 0.706 & 0.725 &  &  & \\ 
%  & \text{Diff.} & 0.039 & 0.015 & 0.028 &  &  & \\ 
\hline
\hline
% \multicolumn{1}{ c }{MFD} & 0.06 & 7 & 21 & 0.06 & 5 & 27 & -\\ 
% \multirow{4}{*}{\multicolumn{1}{ c }{DomainIndep}} & \text{Female} & 0.729 & 0.747 & 0.734 & \multirow{4}{*}{0.010} & \multirow{4}{*}{0.086} & \multirow{4}{*}{0.042} \\ 
%  & \text{Male} & 0.725 & 0.694 & 0.702 &  &  & \\ 
%  & \text{Avg.} & 0.727 & 0.721 & 0.718 &  &  & \\ 
%  & \text{Diff.} & 0.004 & 0.053 & 0.031 &  &  & \\ 
% \hline
% \multicolumn{1}{ c }{DomainIndep} & 0.06 & 7 & 21 & 0.06 & 5 & 27 & -\\ 
{\multirow{4}{*}{FairPrune}} & \text{Female} & 0.776 & 0.711 & 0.734 & \multirow{4}{*}{0.007} & \multirow{4}{*}{0.026} & \multirow{4}{*}{0.014} \\ 
 & \text{Male} & 0.721 & 0.725 & 0.720 &  &  & \\ 
 & \text{Avg. $\uparrow$} & 0.748 & 0.718 & 0.727 &  &  & \\ 
 & \text{Diff. $\downarrow$} & 0.055 & 0.014 & 0.014 &  &  & \\ 
\hline
% \cmidrule(lr){1-1}
% \cmidrule(lr){2-2}
% \cmidrule(lr){3-5}
% \cmidrule(ll){6-8}
% \multicolumn{1}{ c }{FairPrune} & 0.06 & 7 & 21 & 0.06 & 5 & 27 & -\\ 
\multicolumn{1}{ c }{\multirow{4}{*}{ME-FairPrune}} & \text{Female} & 0.770 & 0.723 & 0.742 & \multirow{4}{*}{0.006} & \multirow{4}{*}{\textbf{0.020}} & \multirow{4}{*}{\textbf{0.010}} \\ 
 & \text{Male} & 0.739 & 0.728 & 0.730 &  &  & \\ 
 & \text{Avg. $\uparrow$} & 0.755 & 0.725 & 0.736 &  &  & \\ 
 & \text{Diff. $\downarrow$} & 0.032 & 0.005 & 0.012 &  &  & \\ 
 % \hline
\bottomrule
\end{tabular}
\end{table}


% ver 1
% \begin{table}[ht]
% \centering
% \label{Table:compare_isic}
% \scriptsize
% \begin{tabular}{N N N N N N N N}\toprule
% \multicolumn{1}{N }{\text{}} & \multicolumn{1}{c }{\text{}} &\multicolumn{3}{c }{\text{Accuracy}} & \multicolumn{3}{c }{\text{Fairness}}  \\ 
% \cmidrule(lr){3-5}
% \cmidrule(ll){6-8}
% \multicolumn{1}{ c }{\text{Method}} & \text{Gender} &\text{Precision} & \text{Recall} & \text{F1-score} & \text{Eopp0} & \text{Eopp1} & \text{Eodd} \\ 
% % \cmidrule(lr){1-1}
% % \cmidrule(lr){2-2}
% % \cmidrule(lr){3-5}
% % \cmidrule(ll){6-8}
% \hline
% \multirow{4}{*}{\multicolumn{1}{ c }{ResNet18}} & \text{Female} & 0.793 & 0.721 & 0.746 & \multirow{4}{*}{0.006} & \multirow{4}{*}{0.044} & \multirow{4}{*}{0.022} \\ 
%  & \text{Male} & 0.731 & 0.725 & 0.723 &  &  & \\ 
%  & \text{Avg.} & 0.762 & 0.723 & 0.735 &  &  & \\ 
%  & \text{Diff.} & 0.063 & 0.004 & 0.023 &  &  & \\ 
% \hline

% \multirow{4}{*}{\multicolumn{1}{ c }{AdvConf}} & \text{Female} & 0.755 & 0.738 & 0.741 & \multirow{4}{*}{0.008} & \multirow{4}{*}{0.070} & \multirow{4}{*}{0.037} \\ 
%  & \text{Male} & 0.710 & 0.757 & 0.731 &  &  & \\ 
%  & \text{Avg.} & 0.733 & 0.747 & 0.736 &  &  & \\ 
%  & \text{Diff.} & 0.045 & 0.020 & 0.010 &  &  & \\ 
% % \hline
% \cmidrule(lr){1-1}
% \cmidrule(lr){2-2}
% \cmidrule(lr){3-5}
% \cmidrule(ll){6-8}
% \multirow{4}{*}{\multicolumn{1}{ c }{AdvRev}} & \text{Female} & 0.778 & 0.683 & 0.716 & \multirow{4}{*}{0.007} & \multirow{4}{*}{0.059} & \multirow{4}{*}{0.033} \\ 
%  & \text{Male} & 0.773 & 0.706 & 0.729 &  &  & \\ 
%  & \text{Avg.} & 0.775 & 0.694 & 0.723 &  &  & \\ 
%  & \text{Diff.} & 0.006 & 0.023 & 0.014 &  &  & \\ 
% % \hline
% \multirow{4}{*}{\multicolumn{1}{ c }{DomainIndep}} & \text{Female} & 0.729 & 0.747 & 0.734 & \multirow{4}{*}{0.010} & \multirow{4}{*}{0.086} & \multirow{4}{*}{0.042} \\ 
%  & \text{Male} & 0.725 & 0.694 & 0.702 &  &  & \\ 
%  & \text{Avg.} & 0.727 & 0.721 & 0.718 &  &  & \\ 
%  & \text{Diff.} & 0.004 & 0.053 & 0.031 &  &  & \\ 
% \hline

% \multirow{4}{*}{\multicolumn{1}{ c }{HSIC}} & \text{Female} & 0.744 & 0.660 & 0.696 & \multirow{4}{*}{0.008} & \multirow{4}{*}{0.042} & \multirow{4}{*}{0.019} \\ 
%  & \text{Male} & 0.718 & 0.697 & 0.705 &  &  & \\ 
%  & \text{Avg.} & 0.731 & 0.679 & 0.700 &  &  & \\ 
%  & \text{Diff.} & 0.026 & 0.037 & 0.009 &  &  & \\ 
% \hline

% % \multicolumn{1}{ c }{HSIC} & 0.74 & 15 & 630 & 0.70 & 15 & 618 & -\\ 
% \multirow{4}{*}{\multicolumn{1}{ c }{MFD}} & \text{Female} & 0.770 & 0.697 & 0.727 & \multirow{4}{*}{0.005} & \multirow{4}{*}{0.051} & \multirow{4}{*}{0.024} \\ 
%  & \text{Male} & 0.772 & 0.726 & 0.744 &  &  & \\ 
%  & \text{Avg.} & 0.771 & 0.712 & 0.736 &  &  & \\ 
%  & \text{Diff.} & 0.002 & 0.029 & 0.017 &  &  & \\ 
% \hline

% % \multicolumn{1}{ c }{MFD} & 0.06 & 7 & 21 & 0.06 & 5 & 27 & -\\ 
% % \multirow{4}{*}{\multicolumn{1}{ c }{DomainIndep}} & \text{Female} & 0.729 & 0.747 & 0.734 & \multirow{4}{*}{0.010} & \multirow{4}{*}{0.086} & \multirow{4}{*}{0.042} \\ 
% %  & \text{Male} & 0.725 & 0.694 & 0.702 &  &  & \\ 
% %  & \text{Avg.} & 0.727 & 0.721 & 0.718 &  &  & \\ 
% %  & \text{Diff.} & 0.004 & 0.053 & 0.031 &  &  & \\ 
% % \hline
% % \multicolumn{1}{ c }{DomainIndep} & 0.06 & 7 & 21 & 0.06 & 5 & 27 & -\\ 
% \multirow{4}{*}{\multicolumn{1}{ c }{FairPrune}} & \text{Female} & 0.776 & 0.711 & 0.734 & \multirow{4}{*}{0.007} & \multirow{4}{*}{0.026} & \multirow{4}{*}{0.014} \\ 
%  & \text{Male} & 0.721 & 0.725 & 0.720 &  &  & \\ 
%  & \text{Avg.} & 0.748 & 0.718 & 0.727 &  &  & \\ 
%  & \text{Diff.} & 0.055 & 0.014 & 0.014 &  &  & \\ 
% \hline
% % \multicolumn{1}{ c }{FairPrune} & 0.06 & 7 & 21 & 0.06 & 5 & 27 & -\\ 
% \multirow{4}{*}{\multicolumn{1}{ c }{ME-FairPrune}} & \text{Female} & 0.770 & \textbf{0.723} & 0.742 & \multirow{4}{*}{0.006} & \multirow{4}{*}{\textbf{0.020}} & \multirow{4}{*}{\textbf{0.010}} \\ 
%  & \text{Male} & 0.739 & 0.728 & 0.730 &  &  & \\ 
%  & \text{Avg.} & 0.755 & 0.725 & 0.736 &  &  & \\ 
%  & \text{Diff.} & 0.032 & 0.005 & 0.012 &  &  & \\ 
% \bottomrule
% \end{tabular}
% \end{table}


% \begin{tabular}{SSSSSSSS} 
% \toprule
% \multicolumn{7}{c}{Accuracy} {Fairness}\\
% \cmidrule(r){3-5} \cmidrule(lr){6-8}
%     {Method} & {Skin Tone} & {Precision} & {Recall} & {F1-score} & {Eopp0} & {Eopp1} & {Eodd} \\ \midrule
%     1  & 16.128 & +8.872 & 16.128 & 1.402 & 1.373 & -146.6 & -137.6 \\
%     2  & 3.442  & -2.509 & 3.442  & 0.299 & 0.343 & 133.2  & 152.4  \\
%     3  & 1.826  & -0.363 & 1.826  & 0.159 & 0.119 & 168.5  & -161.1 \\
%     4  & 0.993  & -0.429 & 0.993  & 0.086 & 0.08  & 25.6   & 90     \\ \midrule
%     5  & 1.29   & +0.099 & 1.29   & 0.112 & 0.097 & -175.6 & -114.7 \\
%     6  & 0.483  & -0.183 & 0.483  & 0.042 & 0.063 & 22.3   & 122.5  \\
%     7  & 0.766  & -0.475 & 0.766  & 0.067 & 0.039 & 141.6  & -122   \\
%     8  & 0.624  & +0.365 & 0.624  & 0.054 & 0.04  & -35.7  & 90     \\ \midrule
%     9  & 0.641  & -0.466 & 0.641  & 0.056 & 0.045 & 133.3  & -106.3 \\
%     10 & 0.45   & +0.421 & 0.45   & 0.039 & 0.034 & -69.4  & 110.9  \\
%     11 & 0.598  & -0.597 & 0.598  & 0.052 & 0.025 & 92.3   & -109.3 \\ \bottomrule
% \end{tabular}

% % \begin{table}[h]
\centering
\caption{Accuracy and fairness of classification results across different baselines with and without the ME training framework on the ISIC 2019 dataset.}
\setlength{\tabcolsep}{6pt}
\label{Table:compare_isic_me}
\scriptsize
\begin{tabular}{c c c c c c c c}\toprule
% \hline
\multicolumn{1}{c }{\text{}} & \multicolumn{1}{c }{\text{}} &\multicolumn{3}{c }{\text{Accuracy}} & \multicolumn{3}{c }{\text{Fairness}}  \\ 
\cmidrule(lr){3-5}
\cmidrule(ll){6-8}
\multicolumn{1}{ c }{\text{Method}} & \text{Gender} &\text{Precision} & \text{Recall} & \text{F1-score} & \text{Eopp0 $\downarrow$} & \text{Eopp1 $\downarrow$} & \text{Eodd $\downarrow$} \\ 
\hline
\hline

\multicolumn{1}{ c }{\multirow{4}{*}{ResNet18}} & \text{Female} & 0.793 & 0.721 & 0.746 & \multirow{4}{*}{\textbf{0.006}} & \multirow{4}{*}{0.044} & \multirow{4}{*}{0.022} \\ 
 & \text{Male} & 0.731 & 0.725 & 0.723 &  &  & \\ 
 & \text{Avg. $\uparrow$} & 0.762 & 0.723 & 0.735 &  &  & \\ 
 & \text{Diff. $\downarrow$} & 0.063 & 0.004 & 0.023 &  &  & \\ 
\hline
% \cmidrule(lr){1-1}
% \cmidrule(lr){2-2}
% \cmidrule(lr){3-5}
% \cmidrule(ll){6-8}

\multicolumn{1}{ c }{\multirow{4}{*}{ME-ResNet18}} & \text{Female} & 0.748 & 0.724 & 0.733 & \multirow{4}{*}{\textbf{0.006}} & \multirow{4}{*}{\textbf{0.031}} & \multirow{4}{*}{\textbf{0.016}} \\ 
 & \text{Male} & 0.723 & 0.736 & 0.726 &  &  & \\ 
 & \text{Avg. $\uparrow$} & 0.735 & 0.730 & 0.730 &  &  & \\ 
 & \text{Diff. $\downarrow$} & 0.025 & 0.012 & 0.007 &  &  & \\ 
\hline
\hline

\multicolumn{1}{ c }{\multirow{4}{*}{HSIC}} & \text{Female} & 0.744 & 0.660 & 0.696 & \multirow{4}{*}{0.008} & \multirow{4}{*}{0.042} & \multirow{4}{*}{0.020} \\ 
 & \text{Male} & 0.718 & 0.697 & 0.705 &  &  & \\ 
 & \text{Avg. $\uparrow$} & 0.731 & 0.679 & 0.700 &  &  & \\ 
 & \text{Diff. $\downarrow$} & 0.026 & 0.037 & 0.009 &  &  & \\ 
\hline

\multicolumn{1}{ c }{\multirow{4}{*}{ME-HSIC}} & \text{Female} & 0.733 & 0.707 & 0.716 & \multirow{4}{*}{\textbf{0.007}} & \multirow{4}{*}{{\textbf{0.034}}} & \multirow{4}{*}{{\textbf{0.018}}} \\ 
 & \text{Male} & 0.713 & 0.707 & 0.707 &  &  & \\ 
 & \text{Avg. $\uparrow$} & 0.723 & 0.707 & 0.712 &  &  & \\ 
 & \text{Diff. $\downarrow$} & 0.020 & 0.000 & 0.009 &  &  & \\ 
\hline
\hline

% \multicolumn{1}{ c }{HSIC} & 0.74 & 15 & 630 & 0.70 & 15 & 618 & -\\ 
\multicolumn{1}{ c }{\multirow{4}{*}{MFD}} & \text{Female} & 0.770 & 0.697 & 0.726 & \multirow{4}{*}{\textbf{0.005}} & \multirow{4}{*}{0.051} & \multirow{4}{*}{0.024} \\ 
 & \text{Male} & 0.772 & 0.726 & 0.744 &  &  & \\ 
 & \text{Avg. $\uparrow$} & 0.771 & 0.712 & 0.735 &  &  & \\ 
 & \text{Diff. $\downarrow$} & 0.002 & 0.029 & 0.018 &  &  & \\ 
\hline

\multicolumn{1}{ c }{\multirow{4}{*}{ME-MFD}} & \text{Female} & 0.733 & 0.698 & 0.711 & \multirow{4}{*}{\textbf{0.005}} & \multirow{4}{*}{\textbf{0.024}} & \multirow{4}{*}{\textbf{0.012}} \\ 
 & \text{Male} & 0.772 & 0.713 & 0.739 &  &  & \\ 
 & \text{Avg. $\uparrow$} & 0.752 & 0.706 & 0.725 &  &  & \\ 
 & \text{Diff. $\downarrow$} & 0.039 & 0.015 & 0.028 &  &  & \\ 

% \multicolumn{1}{ c }{MFD} & 0.06 & 7 & 21 & 0.06 & 5 & 27 & -\\ 
% \multirow{4}{*}{\multicolumn{1}{ c }{DomainIndep}} & \text{Female} & 0.729 & 0.747 & 0.734 & \multirow{4}{*}{0.010} & \multirow{4}{*}{0.086} & \multirow{4}{*}{0.042} \\ 
%  & \text{Male} & 0.725 & 0.694 & 0.702 &  &  & \\ 
%  & \text{Avg.} & 0.727 & 0.721 & 0.718 &  &  & \\ 
%  & \text{Diff.} & 0.004 & 0.053 & 0.031 &  &  & \\ 
% \hline
% \multicolumn{1}{ c }{DomainIndep} & 0.06 & 7 & 21 & 0.06 & 5 & 27 & -\\ 
% \hline
\bottomrule
\end{tabular}
\end{table}


% \begin{tabular}{SSSSSSSS} 
% \toprule
% \multicolumn{7}{c}{Accuracy} {Fairness}\\
% \cmidrule(r){3-5} \cmidrule(lr){6-8}
%     {Method} & {Skin Tone} & {Precision} & {Recall} & {F1-score} & {Eopp0} & {Eopp1} & {Eodd} \\ \midrule
%     1  & 16.128 & +8.872 & 16.128 & 1.402 & 1.373 & -146.6 & -137.6 \\
%     2  & 3.442  & -2.509 & 3.442  & 0.299 & 0.343 & 133.2  & 152.4  \\
%     3  & 1.826  & -0.363 & 1.826  & 0.159 & 0.119 & 168.5  & -161.1 \\
%     4  & 0.993  & -0.429 & 0.993  & 0.086 & 0.08  & 25.6   & 90     \\ \midrule
%     5  & 1.29   & +0.099 & 1.29   & 0.112 & 0.097 & -175.6 & -114.7 \\
%     6  & 0.483  & -0.183 & 0.483  & 0.042 & 0.063 & 22.3   & 122.5  \\
%     7  & 0.766  & -0.475 & 0.766  & 0.067 & 0.039 & 141.6  & -122   \\
%     8  & 0.624  & +0.365 & 0.624  & 0.054 & 0.04  & -35.7  & 90     \\ \midrule
%     9  & 0.641  & -0.466 & 0.641  & 0.056 & 0.045 & 133.3  & -106.3 \\
%     10 & 0.45   & +0.421 & 0.45   & 0.039 & 0.034 & -69.4  & 110.9  \\
%     11 & 0.598  & -0.597 & 0.598  & 0.052 & 0.025 & 92.3   & -109.3 \\ \bottomrule
% \end{tabular}

% % \begin{table}[ht]
\centering
\label{Table:compare_merge}
\scriptsize
\begin{tabular}{N N N N N N N N N N N N N N N N}\toprule
\multicolumn{1}{N }{\text{}} & \multicolumn{1}{c }{\text{}} &\multicolumn{3}{c }{\text{Accuracy}} & \multicolumn{3}{c }{\text{Fairness}}  \multicolumn{1}{N }{\text{}} & \multicolumn{1}{c }{\text{}} &\multicolumn{3}{c }{\text{Accuracy}} & \multicolumn{3}{c }{\text{Fairness}}\\ 
\cmidrule(lr){3-5}
\cmidrule(ll){6-8}
\cmidrule(ll){9-9}
\cmidrule(ll){10-10}
\cmidrule(ll){11-13}
\cmidrule(ll){14-16}
\multicolumn{1}{ c }{\text{Method}} & \text{Gender} &\text{Precision} & \text{Recall} & \text{F1-score} & \text{Eopp0} & \text{Eopp1} & \text{Eodd} {\text{Method}} & \text{Skin Tone} &\text{Precision} & \text{Recall} & \text{F1-score} & \text{Eopp0} & \text{Eopp1} & \text{Eodd}\\ 
\cmidrule(lr){1-1}
\cmidrule(lr){2-2}
\cmidrule(lr){3-5}
\cmidrule(ll){6-8}
\cmidrule(ll){9-9}
\cmidrule(ll){10-10}
\cmidrule(ll){11-13}
\cmidrule(ll){14-16}
\multirow{8}{*}{\multicolumn{1}{ c }{ResNet18}} & \text{Female} & 0.793 & 0.721 & 0.746 & \multirow{4}{*}{0.006} & \multirow{4}{*}{0.044} & \multirow{4}{*}{0.022} {\multicolumn{1}{ c }{VGG11}} & \text{Dark} & 0.563 & 0.581 & 0.546 & \multirow{4}{*}{0.0013} & \multirow{4}{*}{0.361} & \multirow{4}{*}{0.182} \\ 
 & \text{Male} & 0.731 & 0.725 & 0.723 &  &  & \text{Light} & 0.482 & 0.495 & 0.473 &  &  & \\ 
 & \text{Avg.} & 0.762 & 0.723 & 0.735 &  &  & \text{Avg.} & 0.523 & 0.538 & 0.510 &  &  &  \\ 
 & \text{Diff.} & 0.063 & 0.004 & 0.023 &  &  & \text{Diff.} & 0.081 & 0.086 & 0.073 &  &  &  \\ 
\hline

\multirow{4}{*}{\multicolumn{1}{ c }{AdvConf}} & \text{Female} & 0.755 & 0.738 & 0.741 & \multirow{4}{*}{0.008} & \multirow{4}{*}{0.070} & \multirow{4}{*}{0.037} \\ 
 & \text{Male} & 0.710 & 0.757 & 0.731 &  &  & \\ 
 & \text{Avg.} & 0.733 & 0.747 & 0.736 &  &  & \\ 
 & \text{Diff.} & 0.045 & 0.020 & 0.010 &  &  & \\ 
\hline
\multirow{4}{*}{\multicolumn{1}{ c }{AdvRev}} & \text{Female} & 0.778 & 0.683 & 0.716 & \multirow{4}{*}{0.007} & \multirow{4}{*}{0.059} & \multirow{4}{*}{0.033} \\ 
 & \text{Male} & 0.773 & 0.706 & 0.729 &  &  & \\ 
 & \text{Avg.} & 0.775 & 0.694 & 0.723 &  &  & \\ 
 & \text{Diff.} & 0.006 & 0.023 & 0.014 &  &  & \\ 
\hline
\multirow{4}{*}{\multicolumn{1}{ c }{DomainIndep}} & \text{Female} & 0.729 & 0.747 & 0.734 & \multirow{4}{*}{0.010} & \multirow{4}{*}{0.086} & \multirow{4}{*}{0.042} \\ 
 & \text{Male} & 0.725 & 0.694 & 0.702 &  &  & \\ 
 & \text{Avg.} & 0.727 & 0.721 & 0.718 &  &  & \\ 
 & \text{Diff.} & 0.004 & 0.053 & 0.031 &  &  & \\ 
\hline

\multirow{4}{*}{\multicolumn{1}{ c }{HSIC}} & \text{Female} & 0.744 & 0.660 & 0.696 & \multirow{4}{*}{0.008} & \multirow{4}{*}{0.042} & \multirow{4}{*}{0.019} \\ 
 & \text{Male} & 0.718 & 0.697 & 0.705 &  &  & \\ 
 & \text{Avg.} & 0.731 & 0.679 & 0.700 &  &  & \\ 
 & \text{Diff.} & 0.026 & 0.037 & 0.009 &  &  & \\ 
\hline

% \multicolumn{1}{ c }{HSIC} & 0.74 & 15 & 630 & 0.70 & 15 & 618 & -\\ 
\multirow{4}{*}{\multicolumn{1}{ c }{MFD}} & \text{Female} & 0.770 & 0.697 & 0.727 & \multirow{4}{*}{0.005} & \multirow{4}{*}{0.051} & \multirow{4}{*}{0.024} \\ 
 & \text{Male} & 0.772 & 0.726 & 0.744 &  &  & \\ 
 & \text{Avg.} & 0.771 & 0.712 & 0.736 &  &  & \\ 
 & \text{Diff.} & 0.002 & 0.029 & 0.017 &  &  & \\ 
\hline

% \multicolumn{1}{ c }{MFD} & 0.06 & 7 & 21 & 0.06 & 5 & 27 & -\\ 
% \multirow{4}{*}{\multicolumn{1}{ c }{DomainIndep}} & \text{Female} & 0.729 & 0.747 & 0.734 & \multirow{4}{*}{0.010} & \multirow{4}{*}{0.086} & \multirow{4}{*}{0.042} \\ 
%  & \text{Male} & 0.725 & 0.694 & 0.702 &  &  & \\ 
%  & \text{Avg.} & 0.727 & 0.721 & 0.718 &  &  & \\ 
%  & \text{Diff.} & 0.004 & 0.053 & 0.031 &  &  & \\ 
% \hline
% \multicolumn{1}{ c }{DomainIndep} & 0.06 & 7 & 21 & 0.06 & 5 & 27 & -\\ 
\multirow{4}{*}{\multicolumn{1}{ c }{FairPrune}} & \text{Female} & 0.776 & 0.711 & 0.734 & \multirow{4}{*}{0.007} & \multirow{4}{*}{0.026} & \multirow{4}{*}{0.014} \\ 
 & \text{Male} & 0.721 & 0.725 & 0.720 &  &  & \\ 
 & \text{Avg.} & 0.748 & 0.718 & 0.727 &  &  & \\ 
 & \text{Diff.} & 0.055 & 0.014 & 0.014 &  &  & \\ 
\hline
% \multicolumn{1}{ c }{FairPrune} & 0.06 & 7 & 21 & 0.06 & 5 & 27 & -\\ 
\multirow{4}{*}{\multicolumn{1}{ c }{ME-FairPrune}} & \text{Female} & 0.770 & \textbf{0.723} & 0.742 & \multirow{4}{*}{0.006} & \multirow{4}{*}{\textbf{0.020}} & \multirow{4}{*}{\textbf{0.010}} \\ 
 & \text{Male} & 0.739 & 0.728 & 0.730 &  &  & \\ 
 & \text{Avg.} & 0.755 & 0.725 & 0.736 &  &  & \\ 
 & \text{Diff.} & 0.032 & 0.005 & 0.012 &  &  & \\ 
\bottomrule
\end{tabular}
\end{table}


% \begin{tabular}{SSSSSSSS} 
% \toprule
% \multicolumn{7}{c}{Accuracy} {Fairness}\\
% \cmidrule(r){3-5} \cmidrule(lr){6-8}
%     {Method} & {Skin Tone} & {Precision} & {Recall} & {F1-score} & {Eopp0} & {Eopp1} & {Eodd} \\ \midrule
%     1  & 16.128 & +8.872 & 16.128 & 1.402 & 1.373 & -146.6 & -137.6 \\
%     2  & 3.442  & -2.509 & 3.442  & 0.299 & 0.343 & 133.2  & 152.4  \\
%     3  & 1.826  & -0.363 & 1.826  & 0.159 & 0.119 & 168.5  & -161.1 \\
%     4  & 0.993  & -0.429 & 0.993  & 0.086 & 0.08  & 25.6   & 90     \\ \midrule
%     5  & 1.29   & +0.099 & 1.29   & 0.112 & 0.097 & -175.6 & -114.7 \\
%     6  & 0.483  & -0.183 & 0.483  & 0.042 & 0.063 & 22.3   & 122.5  \\
%     7  & 0.766  & -0.475 & 0.766  & 0.067 & 0.039 & 141.6  & -122   \\
%     8  & 0.624  & +0.365 & 0.624  & 0.054 & 0.04  & -35.7  & 90     \\ \midrule
%     9  & 0.641  & -0.466 & 0.641  & 0.056 & 0.045 & 133.3  & -106.3 \\
%     10 & 0.45   & +0.421 & 0.45   & 0.039 & 0.034 & -69.4  & 110.9  \\
%     11 & 0.598  & -0.597 & 0.598  & 0.052 & 0.025 & 92.3   & -109.3 \\ \bottomrule
% \end{tabular}

% \textbf{ISIC 2019 Dataset.}
% Table \ref{Table:compare_isic} presents the accuracy and fairness results of all methods and the results of our framework applied on multiple state-of-the-art method on the ISIC 2019 dataset. It can be seen that our ME framework significantly improved the Eodd of the original ResNet18, which considers the difference between true positive rate and false positive rate. This result shows that our framework has the potential to achieve fairness without using sensitive attributes during training. Additionally, our ME framework improves all fairness scores and accuracy in all average scores when applied to FairPrune. We also utilized our ME framework on MFD and HSIC, which initially had better fairness performance than other baselines. With our framework, these models showed better fairness while maintaining similar levels of accuracy. 

% % Table \ref{Table:compare_isic} presents the accuracy and fairness results of all methods on the ISIC 2019 dataset. Our ME framework improves all fairness scores and accuracy in all average scores when applied to FairPrune. We also utilized our ME framework on MFD and HSIC, which initially had better fairness performance than other baselines. With our framework, these models showed better fairness while maintaining similar levels of accuracy. Additionally, our ME framework significantly improved the Eodd of the original ResNet18, which considers the difference between true positive rate and false positive rate. This result shows that our framework has the potential to achieve fairness without using sensitive attributes during training. 


% \subsection{Results on Fitzpatrick-17k Dataset.}

% % In this section, we show the feasibility of ME training framework on three state-of-the-art, HSIC \cite{quadrianto2019discovering}, MFD \cite{jung2021fair} and FairPrune \cite{wu2022fairprune} by comparing the result with and without the ME training framework on Fitzpatrick-17k and ISIC 2019 dataset in Table \ref{Table:Compare_SOTA_CelebA} and Table \ref{Table:Compare_SOTA_UTK}, respectively.

% \begin{table}[h]
\centering
\caption{Results of accuracy and fairness of different methods on Fitzpatrick-17k dataset, using skin tone as the sensitive attribute. The dark skin is the privileged group with higher accuracy by vanilla training. 
The results of \emph{VGG-11}, \emph{AdvConf}, \emph{AdvRev}, \emph{DomainIndep} and \emph{FairPrune} are the experimental results reported in \cite{wu2022fairprune}.
}
\setlength{\tabcolsep}{6pt}
\label{Table:compare_fitz}
\scriptsize
\begin{tabular}{c c c c c c c c}\toprule
% \hline
\multicolumn{1}{c}{\text{}} & \multicolumn{1}{c }{\text{}} &\multicolumn{3}{c }{\text{Accuracy}} & \multicolumn{3}{c }{\text{Fairness}}  \\ 
\cmidrule(lr){3-5}
\cmidrule(ll){6-8}
\multicolumn{1}{ c }{\text{Method}} & \text{Skin Tone} &\text{Precision} & \text{Recall} & \text{F1-score} & \text{Eopp0 $\downarrow$} & \text{Eopp1 $\downarrow$} & \text{Eodd $\downarrow$} \\ 
% \cmidrule(lr){1-1}
% \cmidrule(lr){2-2}
% \cmidrule(lr){3-5}
% \cmidrule(ll){6-8}
\hline
\hline
\multirow{4}{*}{VGG-11} & \text{Dark} & 0.563 & 0.581 & 0.546 & \multirow{4}{*}{0.0013} & \multirow{4}{*}{0.361} & \multirow{4}{*}{0.182} \\ 
 & \text{Light} & 0.482 & 0.495 & 0.473 &  &  & \\ 
 & \text{Avg. $\uparrow$} & 0.523 & 0.538 & 0.510 &  &  & \\ 
 & \text{Diff. $\downarrow$} & 0.081 & 0.086 & 0.073 &  &  & \\ 
\hline
\multicolumn{1}{ c }{\multirow{4}{*}{AdvConf}} & \text{Dark} & 0.506 & 0.562 & 0.506 & \multirow{4}{*}{0.0011} & \multirow{4}{*}{0.339} & \multirow{4}{*}{0.169} \\ 
 & \text{Light} & 0.427 & 0.464 & 0.426 &  &  & \\ 
 & \text{Avg. $\uparrow$} & 0.467 & 0.513 & 0.466 &  &  & \\ 
 & \text{Diff. $\downarrow$} & 0.079 & 0.098 & 0.080 &  &  & \\ 
\hline
\multicolumn{1}{ c }{\multirow{4}{*}{AdvRev}} & \text{Dark} & 0.514 & 0.545 & 0.503 & \multirow{4}{*}{0.0011} & \multirow{4}{*}{0.334} & \multirow{4}{*}{0.166} \\ 
 & \text{Light} & 0.489 & 0.469 & 0.457 &  &  & \\ 
 & \text{Avg. $\uparrow$} & 0.502 & 0.507 & 0.480 &  &  & \\ 
 & \text{Diff. $\downarrow$} & 0.025 & 0.076 & 0.046 &  &  & \\ 
\hline
\multicolumn{1}{ c }{\multirow{4}{*}{DomainIndep}} & \text{Dark} & 0.559 & 0.540 & 0.530 & \multirow{4}{*}{0.0012} & \multirow{4}{*}{0.323} & \multirow{4}{*}{0.161} \\ 
 & \text{Light} & 0.541 & 0.529 & 0.512 &  &  & \\ 
 & \text{Avg. $\uparrow$} & 0.550 & 0.534 & 0.521 &  &  & \\ 
 & \text{Diff. $\downarrow$} & 0.018 & 0.010 & 0.018 &  &  & \\ 
\hline
\multicolumn{1}{ c }{\multirow{4}{*}{HSIC}} & \text{Dark} & 0.548 & 0.522 & 0.513 & \multirow{4}{*}{0.0013} & \multirow{4}{*}{0.331} & \multirow{4}{*}{0.166} \\ 
 & \text{Light} & 0.513 & 0.506 & 0.486 &  &  & \\ 
 & \text{Avg.$\uparrow$} & 0.530 & 0.515 & 0.500 &  &  & \\ 
 & \text{Diff. $\downarrow$} & 0.040 & 0.018 & 0.029 &  &  & \\ 
\hline

% \multicolumn{1}{ c }{HSIC} & 0.74 & 15 & 630 & 0.70 & 15 & 618 & -\\ 
\multicolumn{1}{ c }{\multirow{4}{*}{MFD}} & \text{Dark} & 0.514 & 0.545 & 0.503 & \multirow{4}{*}{0.0011} & \multirow{4}{*}{0.334} & \multirow{4}{*}{0.166} \\ 
 & \text{Light} & 0.489 & 0.469 & 0.457 &  &  & \\ 
 & \text{Avg. $\uparrow$} & 0.502 & 0.507 & 0.480 &  &  & \\ 
 & \text{Diff. $\downarrow$} & 0.025 & 0.076 & 0.046 &  &  & \\ 
\hline
\hline
% \multicolumn{1}{ c }{MFD} & 0.06 & 7 & 21 & 0.06 & 5 & 27 & -\\ 
% \multirow{4}{*}{\multicolumn{1}{ c }{DomainIndep}} & \text{Dark} & 0.559 & 0.540 & 0.530 & \multirow{4}{*}{0.0012} & \multirow{4}{*}{0.323} & \multirow{4}{*}{0.161} \\ 
%  & \text{Light} & 0.541 & 0.529 & 0.512 &  &  & \\ 
%  & \text{Avg.} & 0.550 & 0.534 & 0.521 &  &  & \\ 
%  & \text{Diff.} & 0.018 & 0.010 & 0.018 &  &  & \\ 
% \hline
% \multicolumn{1}{ c }{DomainIndep} & 0.06 & 7 & 21 & 0.06 & 5 & 27 & -\\ 
{\multirow{4}{*}{FairPrune}} & \text{Dark} & 0.567 & 0.519 & 0.507 & \multirow{4}{*}{\textbf{0.0008}} & \multirow{4}{*}{0.330} & \multirow{4}{*}{0.165} \\ 
 & \text{Light} & 0.496 & 0.477 & 0.459 &  &  & \\ 
 & \text{Avg. $\uparrow$} & 0.531 & 0.498 & 0.483 &  &  & \\ 
 & \text{Diff. $\downarrow$} & 0.071 & 0.042 & 0.048 &  &  & \\ 
\hline
% \multicolumn{1}{ c }{FairPrune} & 0.06 & 7 & 21 & 0.06 & 5 & 27 & -\\ 
\multicolumn{1}{ c }{\multirow{4}{*}{ME-FairPrune}} & \text{Dark} & 0.564 & 0.529 & 0.523 & \multirow{4}{*}{0.0012} & \multirow{4}{*}{\textbf{0.305}} & \multirow{4}{*}{\textbf{0.152}} \\ 
 & \text{Light} & 0.542 & 0.535 & 0.522 &  &  & \\ 
 & \text{Avg. $\uparrow$} & 0.553 & 0.532 & 0.522 &  &  & \\ 
 & \text{Diff. $\downarrow$} & 0.022 & 0.006 & 0.001 &  &  & \\
 % \hline
\bottomrule
\end{tabular}
\end{table}


% \begin{tabular}{SSSSSSSS} 
% \toprule
% \multicolumn{7}{c}{Accuracy} {Fairness}\\
% \cmidrule(r){3-5} \cmidrule(lr){6-8}
%     {Method} & {Skin Tone} & {Precision} & {Recall} & {F1-score} & {Eopp0} & {Eopp1} & {Eodd} \\ \midrule
%     1  & 16.128 & +8.872 & 16.128 & 1.402 & 1.373 & -146.6 & -137.6 \\
%     2  & 3.442  & -2.509 & 3.442  & 0.299 & 0.343 & 133.2  & 152.4  \\
%     3  & 1.826  & -0.363 & 1.826  & 0.159 & 0.119 & 168.5  & -161.1 \\
%     4  & 0.993  & -0.429 & 0.993  & 0.086 & 0.08  & 25.6   & 90     \\ \midrule
%     5  & 1.29   & +0.099 & 1.29   & 0.112 & 0.097 & -175.6 & -114.7 \\
%     6  & 0.483  & -0.183 & 0.483  & 0.042 & 0.063 & 22.3   & 122.5  \\
%     7  & 0.766  & -0.475 & 0.766  & 0.067 & 0.039 & 141.6  & -122   \\
%     8  & 0.624  & +0.365 & 0.624  & 0.054 & 0.04  & -35.7  & 90     \\ \midrule
%     9  & 0.641  & -0.466 & 0.641  & 0.056 & 0.045 & 133.3  & -106.3 \\
%     10 & 0.45   & +0.421 & 0.45   & 0.039 & 0.034 & -69.4  & 110.9  \\
%     11 & 0.598  & -0.597 & 0.598  & 0.052 & 0.025 & 92.3   & -109.3 \\ \bottomrule
% \end{tabular}

% \textbf{Fitzpatrick-17k Dataset.}
% To evaluate the extensibility of our framework on different model structures, we utilized VGG11 as the backbone on the Fitzpatrick-17k dataset. Table \ref{Table:compare_fitz} presents the accuracy and fairness results of all methods, demonstrating that the ME framework outperformed all other methods with the best EOpp1 and EOdd scores, showing a 7.5\% and 7.9\% improvement over FairPrune. Additionally, the F1-score difference and average value were superior compared to other methods. Our results also showed better mean accuracy and smaller accuracy differences in all criteria compared to FairPrune. This is because the ME framework ensures that early exit instances have a high level of confidence in their correctness, and the classification through shallower, fairer features further improves fairness.

% \subsection{Ablation Study of Multi-exit Training Framework}

% \begin{figure*}[ht]
% % \setlength{\belowcaptionskip}{-20pt} 
% \begin{center}

% \includegraphics[width=0.99\linewidth]{Fig/Ablation.png}
% % \includegraphics[width=0.9\columnwidth]{Fig/Method.png}
% \caption{Ablation study for our paper.
% }
% \label{Fig.Ablation}
% \end{center}
% \end{figure*}

% In this section, we conduct an ablation study to investigate the effectiveness of applying a multi-exit training framework to the ME-FairPrune method \cite{wu2022fairprune}. We apply ME-FairPrune to a pre-trained model from the ISIC 2019 dataset and vary the confidence threshold to examine its impact on accuracy and fairness. Our results in fig \ref{Fig.Ablation} (a) show that increasing the threshold improves both accuracy and fairness. Moreover, when the threshold exceeds 0.9, the fairness score outperforms the original FairPrune method. We recommend setting the threshold based on validation set results to ensure optimal performance.   
% In fig \ref{Fig.Ablation} (b) and (c), we compare the fairness and accuracy between ME-FairPrune using an early exit policy and exiting from a specific exit. Our findings in fig \ref{Fig.Ablation}(b) indicate that, across all criteria, no $IC$ outperforms the early exit policy in terms of fairness metrics, while the early exit policy achieves accuracy levels comparable to those of $IC5$ with better EOdd and EOpp1. As for \ref{Fig.Ablation}(c) shows that with early exit policy, model can keep similar accuracy level compare to results inference from certain IC. 


% Table \ref{Table:compare_isic} presents the accuracy and fairness results of all methods on the ISIC 2019 dataset. Our ME framework improves all fairness scores and accuracy in all average scores when applied to FairPrune. We also utilized our ME framework on MFD and HSIC, which initially had better fairness performance than other baselines. With our framework, these models showed better fairness while maintaining similar levels of accuracy. Additionally, our ME framework significantly improved the Eodd of the original ResNet18, which considers the difference between true positive rate and false positive rate. This result shows that our framework has the potential to achieve fairness without using sensitive attributes during training. 


% \subsection{Results on Fitzpatrick-17k Dataset.}

% In this section, we show the feasibility of ME training framework on three state-of-the-art, HSIC \cite{quadrianto2019discovering}, MFD \cite{jung2021fair} and FairPrune \cite{wu2022fairprune} by comparing the result with and without the ME training framework on Fitzpatrick-17k and ISIC 2019 dataset in Table \ref{Table:Compare_SOTA_CelebA} and Table \ref{Table:Compare_SOTA_UTK}, respectively.